\chapter{Games, Label Cover, and the Language of PCPs}
\label{lec:games}

{\sffamily\large\color{softgray}One object with four names, and the price of every translation\par}
\vspace{0.5em}

\noindent
Lecture~\ref{lec:why} introduced four characters: a PCP verifier, a constraint satisfaction
problem, a two-prover game, and a Label Cover instance.  Table~\ref{tab:dict} asserted that they are
one and the same object seen from four sides.  That is true, but only up to constants, and in hardness of
approximation constants are exactly what we cannot afford to lose track of: the difference between
soundness $1-\eps_0$ and soundness $1-\eps_0/3$ is invisible to the PCP theorem and decisive for a reduction
that must end with soundness $\tfrac12+\eps$.

This lecture turns every arrow of the dictionary into a theorem and records what it costs.  Three prices
will appear: a factor $2$ in value when we pass from one oracle to two provers (Section~\ref{sec:oracle}), a
factor $q$ in the gap when we pass from $q$-ary to binary constraints (Section~\ref{sec:cv}), and a blow-up
of the alphabet.  Once the prices are known, statements such as ``Label Cover with soundness $\delta$ over
an alphabet of size $\poly(1/\delta)$'' or ``$2$-to-$2$ Games over $\F_2^\ell$'' can be read at a glance.
We end with the oldest customer of this language, the reduction of
\citet{FGLSS96} from PCPs to \sat{Max-Clique}.  Its proof fits on one page and contains the idea behind
every later reduction in this course: \emph{turn the consistency of local views of a proof into the structure of a graph}.

\begin{goalbox}{After this lecture you should be able to\ldots}
\begin{enumerate}[leftmargin=*,label=(\arabic*)]
  \item write any PCP verifier, CSP, game or Label Cover instance as a constraint graph, and compute
        the value of a small one by hand (Sections~\ref{sec:lang}, \ref{sec:cv});
  \item say which constraint types---general, projection, $2$-to-$2$, unique---a given instance has, and
        what each type buys and costs (Section~\ref{sec:ladder});
  \item convert between a one-oracle and a two-prover view of the same instance, and state exactly how
        much value is lost (Proposition~\ref{prop:cover});
  \item prove that the clause--variable game of any $q$-ary CSP has soundness $1-\eps/q$ (Theorem~\ref{thm:cv});
  \item build the FGLSS graph of a verifier and prove that its clique number is $2^r$ times the best
        acceptance probability (Theorem~\ref{thm:fglss}).
\end{enumerate}
\end{goalbox}


%% ===========================================================================
\section{Constraint graphs: one language for everything}
\label{sec:lang}
%% ===========================================================================

We begin with the most general object and specialize.

\begin{definition}[Constraint satisfaction instance]\label{def:csp}
An instance $\Phi$ of a \defn{constraint satisfaction problem} consists of a finite set $X$ of
\emph{variables}, a finite alphabet $\Sigma_x$ for every $x\in X$, and a finite multiset $\mathcal C$ of
\emph{constraints}.  A constraint $c=(S_c,P_c)$ has a \emph{scope} $S_c\subseteq X$ and a
\emph{predicate} $P_c\colon\prod_{x\in S_c}\Sigma_x\to\{0,1\}$.  The \emph{arity} of $\Phi$ is $q=\max_c|S_c|$.
An \emph{assignment} $\sigma$ gives every variable $x$ a value $\sigma(x)\in\Sigma_x$, and
\[
  \val(\Phi)\;=\;\max_{\sigma}\ \Pr_{c\sim\mathcal C}\bigl[P_c(\sigma|_{S_c})=1\bigr].
\]
\end{definition}

\noindent
(Lecture~\ref{lec:why} wrote $\OPT$ for the optimum of a Max-CSP; from now on we write $\val$, which is also
the name for the value of a game.)  The problems of Lecture~\ref{lec:why} all fit.

\begin{table}[t]
\centering\small
\renewcommand{\arraystretch}{1.3}
\begin{tabular}{@{}>{\raggedright\arraybackslash}p{3.3cm}
                   >{\raggedright\arraybackslash}p{2.1cm}
                   c
                   >{\raggedright\arraybackslash}p{3.7cm}
                   >{\raggedright\arraybackslash}p{2.2cm}@{}}
\toprule
Problem & Alphabet & Arity & Predicate & Random assignment\\
\midrule
\sat{Max-E3-SAT} & $\{0,1\}$ & $3$ & OR of three literals & $\tfrac78$\\
\sat{Max-E3-Lin}$(\F_2)$ & $\{0,1\}$ & $3$ & $x+y+z=b$ & $\tfrac12$\\
\sat{Max-Cut} & $\{0,1\}$ & $2$ & $x\ne y$ & $\tfrac12$\\
\sat{Max-3-Cut} & $\{0,1,2\}$ & $2$ & $x\ne y$ & $\tfrac23$\\
Linear Unique Games & $\Z_k$ & $2$ & $x-y=c$ & $\tfrac1k$\\
Label Cover (projection) & $\Sigma_U,\Sigma_W$ & $2$ & $b=\pi(a)$ & $\tfrac1{|\Sigma_W|}$\\
$2$-to-$2$ Games & $[k]$ & $2$ & $\pi(a)=\pi'(b)$ & $\tfrac2k$\\
\bottomrule
\end{tabular}
\caption{A zoo of constraint satisfaction problems.  The last column is the fraction of constraints satisfied
by a \emph{uniformly random assignment}; it is the ``trivial'' threshold that approximation algorithms must beat.}
\label{tab:zoo}
\end{table}

The last column of Table~\ref{tab:zoo} is not a coincidence of the examples.  It is an expectation.

\begin{lemma}[Random assignment]\label{lem:random-asg}
Let the \defn{density} of a constraint $c$ be $\delta(c)=|P_c^{-1}(1)|\big/\prod_{x\in S_c}|\Sigma_x|$.  Then
$\val(\Phi)\ge\E_{c\sim\mathcal C}\,\delta(c)$.
\end{lemma}
\begin{proof}
For a uniformly random assignment $\sigma$ and a fixed constraint $c$, the restriction $\sigma|_{S_c}$ is uniform on
$\prod_{x\in S_c}\Sigma_x$, so $c$ is satisfied with probability $\delta(c)$.  By linearity of expectation
the expected fraction of satisfied constraints is $\E_c\,\delta(c)$, and some assignment does at least as well.
\end{proof}

\noindent
This is the baseline of the scoreboard in Table~\ref{tab:scoreboard}: the first two rows say that for
\sat{Max-E3-SAT} and \sat{Max-E3-Lin} nothing efficient beats the density of the predicate, unless $\Ptime=\NP$.

\subsection{Binary constraints: constraint graphs}

From here on the arity is almost always $2$.

\begin{definition}[Constraint graph]\label{def:cg}
A \defn{constraint graph} is a CSP instance of arity $2$.  Concretely it is a finite multigraph $G=(V,E)$,
alphabets $\Sigma_v$ for $v\in V$, and, for every edge, oriented arbitrarily as $e=(u,v)$, a
\defn{relation} $C_e\subseteq\Sigma_u\times\Sigma_v$.  A \emph{labeling} $\ell$ gives each vertex $v$ a label
$\ell(v)\in\Sigma_v$; it \emph{satisfies} $e=(u,v)$ if $(\ell(u),\ell(v))\in C_e$, and
\[
  \val(G)\;=\;\max_{\ell}\ \Pr_{e=(u,v)\sim E}\bigl[(\ell(u),\ell(v))\in C_e\bigr].
\]
If $G$ is bipartite, $G=(U,W,E)$ with every edge oriented from $U$ to $W$, we call it a \defn{game}; this is
Definition~\ref{def:game}, and $\val(G)$ is the value of the game.
\end{definition}

\noindent
When an edge $e=(u,v)$ carries a projection we write $\pi_{uv}\colon\Sigma_u\to\Sigma_v$ so that
$C_e=\{(a,\pi_{uv}(a))\}$, and the value reads
$\val(G)=\max_\ell\Pr_{(u,v)\sim E}[\pi_{uv}(\ell(u))=\ell(v)]$, the formula we shall use throughout.
Unique Games (Definition~\ref{def:ug}) are the case where every $\pi_{uv}$ is a bijection of $[k]$.

The word ``Label Cover'' is attached to this optimization problem---find a labeling of maximum value---and
usually, though not always, refers to projection constraints.  We shall say \emph{Label Cover} for the
general problem and specify the constraint type when it matters.

\begin{definition}[Types of constraints]\label{def:types}
Let $C\subseteq\Sigma\times\Sigma'$ be the relation on an edge $(u,v)$, with $\Sigma=\Sigma_u$ and $\Sigma'=\Sigma_v$.
\begin{itemize}
  \item $C$ is a \defn{projection} if $C=\{(a,\pi(a))\}$ for a function $\pi\colon\Sigma\to\Sigma'$: every left label has
        exactly one partner.  It is \defn{$d$-to-$1$} if moreover every fibre $\pi^{-1}(b)$ has exactly $d$ elements.
  \item $C$ is \defn{unique} if it is the graph of a bijection, \ie a perfect matching between $\Sigma$ and $\Sigma'$.
  \item $C$ is \defn{$2$-to-$2$} if there are a set $\Gamma$ and two $2$-to-$1$ maps $\pi\colon\Sigma\to\Gamma$,
        $\pi'\colon\Sigma'\to\Gamma$ with $C=\{(a,b):\pi(a)=\pi'(b)\}$.  Every label then has exactly two partners,
        and the bipartite graph of $C$ is a disjoint union of copies of $K_{2,2}$.
\end{itemize}
An instance \emph{is} of a given type if all its constraints are.
\end{definition}

\begin{figure}[t]
\centering
\begin{tikzpicture}[x=1cm,y=1cm,every node/.style={font=\small},
   la/.style={circle,fill=accent,inner sep=1.6pt}, lb/.style={circle,fill=warn,inner sep=1.6pt},
   ed/.style={thick,black!70}]
  % (a) arbitrary relation
  \begin{scope}
    \node[font=\sffamily\footnotesize\bfseries] at (0.95,0.75) {arbitrary};
    \foreach \i in {1,...,4} {\node[la] (pa\i) at (0,-\i*0.6+0.6) {}; \node[lb] (pb\i) at (1.9,-\i*0.6+0.6) {};}
    \draw[ed] (pa1)--(pb1) (pa1)--(pb2) (pa1)--(pb3) (pa2)--(pb2) (pa3)--(pb1) (pa3)--(pb4);
    \node[font=\footnotesize] at (0.95,-2.3) {general relation};
  \end{scope}
  % (b) projection 2-to-1
  \begin{scope}[xshift=4.0cm]
    \node[font=\sffamily\footnotesize\bfseries] at (0.95,0.75) {projection};
    \foreach \i in {1,...,4} \node[la] (qa\i) at (0,-\i*0.6+0.6) {};
    \node[lb] (qb1) at (1.9,-0.3) {}; \node[lb] (qb2) at (1.9,-1.5) {};
    \draw[ed] (qa1)--(qb1) (qa2)--(qb1) (qa3)--(qb2) (qa4)--(qb2);
    \node[font=\footnotesize] at (0.95,-2.3) {here $2$-to-$1$};
  \end{scope}
  % (c) 2-to-2: two matchings
  \begin{scope}[xshift=8.0cm]
    \node[font=\sffamily\footnotesize\bfseries] at (0.95,0.75) {2-to-2};
    \foreach \i in {1,...,4} {\node[la] (ra\i) at (0,-\i*0.6+0.6) {}; \node[lb] (rb\i) at (1.9,-\i*0.6+0.6) {};}
    \draw[accent,thick] (ra1)--(rb1) (ra2)--(rb2) (ra3)--(rb3) (ra4)--(rb4);
    \draw[warn,thick,dashed] (ra1)--(rb2) (ra2)--(rb1) (ra3)--(rb4) (ra4)--(rb3);
    \node[font=\footnotesize] at (0.95,-2.3) {two matchings in one};
  \end{scope}
  % (d) unique
  \begin{scope}[xshift=12.0cm]
    \node[font=\sffamily\footnotesize\bfseries] at (0.95,0.75) {unique};
    \foreach \i in {1,...,4} {\node[la] (sa\i) at (0,-\i*0.6+0.6) {}; \node[lb] (sb\i) at (1.9,-\i*0.6+0.6) {};}
    \draw[ed] (sa1)--(sb3) (sa2)--(sb1) (sa3)--(sb4) (sa4)--(sb2);
    \node[font=\footnotesize] at (0.95,-2.3) {perfect matching};
  \end{scope}
\end{tikzpicture}
\caption{The four constraint types of Definition~\ref{def:types} as bipartite graphs between $\Sigma_u$
(left, blue) and $\Sigma_v$ (right, orange).  In a unique constraint each label determines its partner
\emph{and} vice versa; in a projection, only one direction is determined; in a $2$-to-$2$ constraint each label
leaves two candidates.  The $2$-to-$2$ relation of panel~(c) is the union of the solid and the dashed
perfect matching (Lemma~\ref{lem:split}).}
\label{fig:types}
\end{figure}

Lemma~\ref{lem:random-asg} immediately quantifies how many labels a type forces on us.

\begin{corollary}[Alphabets cannot be small]\label{cor:alph}
For a projection game, $\val(G)\ge1/|\Sigma_W|$; for a Unique Game with $k$ labels, $\val(G)\ge1/k$;
for a $2$-to-$2$ game with $k$ labels on both sides, $\val(G)\ge 2/k$.  Hence a hardness result with soundness $\delta$ needs
alphabets of size at least $1/\delta$, $1/\delta$ and $2/\delta$, respectively.
\end{corollary}
\begin{proof}
A projection has $|C|=|\Sigma|$ pairs, hence density $|\Sigma|/(|\Sigma||\Sigma'|)=1/|\Sigma'|$; a $2$-to-$2$ constraint has
$|C|=2|\Sigma|$ pairs and, when $|\Sigma|=|\Sigma'|=k$, density $2/k$.  Apply Lemma~\ref{lem:random-asg}.
\end{proof}

\begin{example}[Unique Games on a cycle]\label{ex:cycle}
Let $G$ be a cycle on vertices $0,1,\dots,\ell-1$ with the linear constraints $x_{i+1}-x_i=c_i\pmod k$ on the edges
$(i,i+1)$ (indices modulo $\ell$).  If $\sum_ic_i\equiv0$, set $x_0=0$ and propagate: all $\ell$ equations hold.
Otherwise summing all the equations would give $0\equiv\sum_ic_i\not\equiv0$, so at least one equation fails; but
deleting one edge leaves a path, which can always be satisfied.  Hence $\val(G)=1-\tfrac1\ell$.
\emph{A cycle is satisfiable iff the constraints compose to the identity around it}; this is the whole
obstruction to perfect completeness in a Unique Game, and checking it along a spanning tree is what the algorithm
of Proposition~\ref{prop:ug-easy} does.
\end{example}


%% ===========================================================================
\section{One oracle or two provers?}
\label{sec:oracle}
%% ===========================================================================

A PCP verifier that makes two queries reads two positions of \emph{one} proof.  In a game the referee sends two
questions to \emph{two} provers who answer independently.  A constraint graph on a general graph is the first,
a game the second.  The two differ in one respect: a single proof is automatically consistent with itself, whereas
two provers may give contradictory answers to the same question.

\begin{definition}[Double cover]\label{def:cover}
Let $G=(V,E)$ be a constraint graph with relations $C_e$.  Its \defn{double cover} $G^{\times2}$ is the game with left
vertices $v_A$ (Alice's questions) and right vertices $v_B$ (Bob's questions), one pair $(v_A,v_B)$ for
every $v\in V$, both with alphabet $\Sigma_v$, and for every edge $e=(u,v)$ of $G$ the two edges
$(u_A,v_B)$ with relation $C_e$ and $(v_A,u_B)$ with the transposed relation $C_e^{\top}=\{(b,a):(a,b)\in C_e\}$.
In words: the referee picks an edge of $G$ and a random orientation of it, and asks the two endpoints to the two provers.
\end{definition}

\begin{proposition}[Value of the double cover]\label{prop:cover}
For every constraint graph $G$,
\[
  \val(G)\;\le\;\val(G^{\times2})\;\le\;2\,\val(G),
\]
and both inequalities can be tight.
\end{proposition}
\begin{proof}
For the first inequality let $\ell$ be an optimal labeling of $G$ and let Alice and Bob both use it:
$\ell_A(v_A)=\ell_B(v_B)=\ell(v)$.  Whatever orientation the referee picks, the two answers are the labels of the two endpoints, so
the acceptance probability is $\val(G)$.

For the second, let $(\ell_A,\ell_B)$ be an optimal strategy in $G^{\times2}$ with value $p$.  Define a labeling of
$G$ at random: for every vertex $v$, independently and uniformly pick a side $s(v)\in\{A,B\}$ and set
$\ell(v)=\ell_{s(v)}(v)$.  Fix an edge $e=(u,v)$.  With probability $\tfrac14$ we have $s(u)=A$ and $s(v)=B$;
then $\ell$ satisfies $e$ iff the edge $(u_A,v_B)$ of $G^{\times2}$ is satisfied by $(\ell_A,\ell_B)$.  With probability
$\tfrac14$ we have $s(v)=A$ and $s(u)=B$, and $\ell$ satisfies $e$ iff the edge $(v_A,u_B)$ is satisfied.  These two events are
disjoint, so
\[
  \Pr[\ell\text{ satisfies }e]\ \ge\ \tfrac14\bigl(\mathbf 1[(u_A,v_B)\text{ sat.}]+\mathbf 1[(v_A,u_B)\text{ sat.}]\bigr).
\]
Averaging over a uniform edge $e$ of $G$, the right-hand side has expectation $\tfrac14\cdot2p=p/2$, since the $2|E|$ edges of
$G^{\times2}$ are the two orientations of the edges of $G$.  Hence some labeling has value at least $p/2$, \ie
$\val(G)\ge\val(G^{\times2})/2$.

For tightness of the first inequality take any satisfiable $G$.  For the second take \sat{Max-Cut} on the complete graph $K_n$:
the best cut contains a $\tfrac{\lceil n/2\rceil\lfloor n/2\rfloor}{\binom n2}\to\tfrac12$ fraction of the edges, but in the double
cover Alice answers $0$ everywhere and Bob $1$ everywhere, and all constraints $a\ne b$ hold (Figure~\ref{fig:cover}).
\end{proof}

\begin{figure}[t]
\centering
\begin{tikzpicture}[x=1cm,y=1cm,every node/.style={font=\small},
   al/.style={circle,draw,fill=accent!22,minimum size=0.75cm,inner sep=0pt},
   bo/.style={circle,draw,fill=warn!22,minimum size=0.75cm,inner sep=0pt},
   pl/.style={circle,draw,fill=black!6,minimum size=0.7cm,inner sep=0pt}]
  % triangle
  \node[pl] (u) at (0,1.5) {$u$};
  \node[pl] (v) at (-1.3,-0.5) {$v$};
  \node[pl] (w) at (1.3,-0.5) {$w$};
  \draw[thick] (u)--(v)--(w)--(u);
  \node[align=center,font=\footnotesize] at (0,-1.7) {\sat{Max-Cut} on a triangle\\ $\val=\tfrac23$};
  \draw[-{Stealth[length=2.2mm]},thick] (2.5,0.5)--(4.1,0.5) node[midway,above,font=\footnotesize] {double cover};
  % hexagon
  \begin{scope}[xshift=7.4cm,yshift=0.35cm]
    \node[al] (uA) at (90:1.55) {$u_A$};
    \node[bo] (vB) at (30:1.55) {$v_B$};
    \node[al] (wA) at (-30:1.55) {$w_A$};
    \node[bo] (uB) at (-90:1.55) {$u_B$};
    \node[al] (vA) at (-150:1.55) {$v_A$};
    \node[bo] (wB) at (150:1.55) {$w_B$};
    \draw[thick] (uA)--(vB)--(wA)--(uB)--(vA)--(wB)--(uA);
    \node[align=center,font=\footnotesize] at (0,-2.05) {the induced game\\ $\val=1$: Alice says $0$, Bob says $1$};
  \end{scope}
\end{tikzpicture}
\caption{The double cover of a triangle is a hexagon, which is bipartite.  Self-consistency of a single labeling is what
the triangle tests; the two provers of the hexagon need not be consistent with each other, and they win with probability~$1$.}
\label{fig:cover}
\end{figure}

\begin{remark}
The factor~$2$ is harmless when the soundness we care about is small: a hardness result for constraint graphs with completeness $c$ and soundness $s$
gives one for games with completeness $c$ and soundness $2s$.  It is fatal when $s\ge\tfrac12$.  This is why the PCP theorem is converted into a game by the
clause--variable construction of Section~\ref{sec:cv} (soundness $1-\eps/q$) and \emph{not} by the double cover
(soundness $2(1-\eps)$, which exceeds $1$ and says nothing as soon as $\eps<\tfrac12$).  For fixed orientations the loss is a factor $4$ instead of~$2$, which is the
bound quoted in \citet{Khot10}.
\end{remark}

A game is conversely a perfectly good PCP.

\begin{lemma}[Answer tables]\label{lem:tables}
Let $G=(U,W,E)$ be a game with alphabets $\Sigma_U,\Sigma_W$ of size at most $\sigma$.  Consider the verifier $V_G$ whose proof is a
pair of tables $A\colon U\to\Sigma_U$ and $B\colon W\to\Sigma_W$, that picks a uniformly random edge $(u,w)$, reads $A(u)$ and $B(w)$, and accepts
iff $(A(u),B(w))\in C_{(u,w)}$.  Then $\max_{A,B}\Pr[V_G\text{ accepts}]=\val(G)$.  If proofs must be binary strings, each label
is written with $\lceil\log_2\sigma\rceil$ bits, and the verifier makes $2\lceil\log_2\sigma\rceil$ binary queries, uses $\log_2|E|$ random bits
and has the same completeness and soundness.
\end{lemma}
\begin{proof}
A proof $(A,B)$ is exactly a labeling of $G$, and the acceptance probability is the fraction of satisfied edges.
\end{proof}

\noindent
So $\Gap{1}{\delta}{LabelCover}_\sigma$ is a PCP with $O(\log\sigma)$ queries and soundness~$\delta$.  Is that better than what we
could already do?  By Lemma~\ref{lem:tables} and Theorem~\ref{thm:lc} the parameters are $O(\log(1/\delta))$ queries and
soundness~$\delta$ with $O(\log n)$ random bits---and independent repetition of the plain verifier of Proposition~\ref{prop:pcp-equiv}
achieves the very same trade-off.  The point of games is not their parameters but their \emph{structure}: the $2$ queries
come in two groups, and the answer in one group is a function of the answer in the other.  This is the structure that the long-code
tests of Lectures~7--11 exploit.

\begin{proposition}[Two binary queries are not enough]\label{prop:2sat}
If a language $L$ has a PCP verifier with $r=O(\log n)$ random bits, \emph{two} binary queries, perfect completeness and soundness $s<1$,
then $L\in\Ptime$.  Hence, unless $\Ptime=\NP$, a two-query PCP for $\NP$ with perfect completeness needs alphabets of growing size.
\end{proposition}
\begin{proof}
On input $x$ enumerate the $2^r=\poly(n)$ random strings $\rho$.  Each defines a predicate $P_\rho$ on two Boolean proof bits,
and every Boolean predicate on two variables is a $2$-CNF (forbid each rejected pair of values by one clause of width $2$).
So the verifier defines a $2$-CNF $\psi_x$ of polynomial size.  If $x\in L$ the honest proof satisfies $\psi_x$; if $x\notin L$ every
proof is rejected for some $\rho$, so $\psi_x$ is unsatisfiable.  $2$SAT is in $\Ptime$.
\end{proof}


%% ===========================================================================
\section{From verifiers and CSPs to games}
\label{sec:cv}
%% ===========================================================================

Lecture~\ref{lec:why} turned the clauses of a $3$-CNF into a game by asking one prover about a clause and the other
about a variable of it.  The construction works for every CSP, and its analysis is the same one-paragraph argument.

\begin{definition}[Constraint--variable game]\label{def:cvgame}
Let $\Phi$ be a CSP instance of arity $q$.  The \defn{constraint--variable game} $G_\Phi$ is the following game.  The referee draws a
uniformly random constraint $c$ and a uniformly random variable $x\in S_c$, and sends $c$ to Alice and $x$ to Bob.  Alice
answers an assignment $a\in\prod_{y\in S_c}\Sigma_y$ to the scope of $c$, Bob answers a symbol $b\in\Sigma_x$, and they win iff
$P_c(a)=1$ and $a_x=b$.
\end{definition}

\noindent
(The edge $(c,x)$ has weight $\tfrac1{|\mathcal C||S_c|}$; weights are harmless since we may repeat edges.  If every predicate is satisfiable we restrict
Alice's alphabet to the satisfying assignments of $c$, which changes nothing since a non-satisfying answer loses outright.  Then $G_\Phi$ is a \emph{projection game}:
$\pi_{(c,x)}(a)=a_x$, with $|\Sigma_U|\le\max_x|\Sigma_x|^q$ and $\Sigma_W=\Sigma_x$.)

\begin{theorem}[Soundness of the constraint--variable game]\label{thm:cv}
For every CSP instance $\Phi$ of arity $q$,
\[
  \val(\Phi)\ \le\ \val(G_\Phi)\ \le\ 1-\frac{1-\val(\Phi)}{q}.
\]
In particular $\val(\Phi)=1$ iff $\val(G_\Phi)=1$, and if every assignment violates an $\eps$-fraction of the constraints, then
$\val(G_\Phi)\le1-\eps/q$.
\end{theorem}
\begin{proof}
\emph{Lower bound.}  Let $\sigma$ be an optimal assignment of $\Phi$.  Bob answers $\sigma(x)$; Alice, for a constraint $c$ that $\sigma$
satisfies, answers $\sigma|_{S_c}$, and for a constraint it violates, anything.  On every constraint satisfied by $\sigma$ the provers win
for every $x\in S_c$, hence the winning probability is at least $\val(\Phi)$.

\emph{Upper bound.}  Fix deterministic strategies.  Bob's answers define an assignment $b$ of $\Phi$ with $b(x)=\ell_W(x)$.
Let $c$ be a constraint violated by $b$, and let $a=\ell_U(c)$ be Alice's answer.  If $P_c(a)=0$ the provers lose on every
edge at~$c$.  Otherwise $a\ne b|_{S_c}$ (because $P_c(b|_{S_c})=0$), so there is some $x\in S_c$ with $a_x\ne b(x)$, and the referee
asks that $x$ with probability at least $1/|S_c|\ge1/q$.  So the provers lose at $c$ with probability at least $1/q$.  Since $b$ violates at
least a $1-\val(\Phi)$ fraction of the constraints, the total losing probability is at least $(1-\val(\Phi))/q$.
\end{proof}

\noindent
This is Proposition~\ref{prop:cv} for $q=3$ and a CNF.  The bound is attained, \eg for the formula below.

\begin{example}[A formula and its game]\label{ex:phi0}
Let $\varphi_0$ have the four clauses $C_1=(x_1\vee x_2)$, $C_2=(\neg x_1\vee x_2)$, $C_3=(\neg x_2\vee x_3)$, $C_4=(\neg x_2\vee\neg x_3)$.
If $x_2=1$ then $C_1,C_2$ hold and exactly one of $C_3,C_4$ does; if $x_2=0$ then $C_3,C_4$ hold and exactly one of $C_1,C_2$ does.
So \emph{every} assignment satisfies exactly three of the four clauses: $\val(\varphi_0)=\tfrac34$.  Theorem~\ref{thm:cv} with $q=2$ gives
$\val(G_{\varphi_0})\le1-\tfrac{1/4}{2}=\tfrac78$.  This is tight: let Bob answer any assignment $b$, say $b=(0,0,0)$, which violates $C_1$ only; Alice answers $b|_{S_c}$ on $C_2,C_3,C_4$ and
$(x_1,x_2)=(1,0)$ on $C_1$, a satisfying assignment that still agrees with $b$ on $x_2$.  The provers lose only if the referee picks $C_1$ and then the variable $x_1$, an event of
probability $\tfrac14\cdot\tfrac12=\tfrac18$.
\end{example}

\begin{corollary}[Every verifier is a game]\label{cor:verifier-game}
Let $V$ be a PCP verifier with $r$ random bits and $q$ queries, perfect completeness and soundness~$s$ on a language~$L$.  Let $\Phi_V(x)$ be the CSP with one
constraint per random string~$\rho$ (scope: the positions queried at~$\rho$, predicate: $P_{x,\rho}$).  Then the game
$G_{\Phi_V(x)}$, in which Alice is asked a random string $\rho$ and Bob a position queried at $\rho$, is a projection game with $|\Sigma_U|\le2^q$, $|\Sigma_W|=2$, and
\[
  x\in L\ \Rightarrow\ \val=1,\qquad x\notin L\ \Rightarrow\ \val\le1-\frac{1-s}{q}.
\]
\end{corollary}
\begin{proof}
$\val(\Phi_V(x))$ is the best acceptance probability over proofs, which is $1$ or at most $s$; apply Theorem~\ref{thm:cv}.
\end{proof}

\noindent
Combined with $\NP=\PCP[O(\log n),O(1)]$, Corollary~\ref{cor:verifier-game} says that it is $\NP$-hard to distinguish projection games of value $1$ from games of value at most $1-\eps_1$, for a constant $\eps_1>0$:
the PCP theorem in the language of games, stated without mentioning $3$SAT.

Table~\ref{tab:price} collects the translations between the four columns of the dictionary and what each costs.

\begin{table}[t]
\centering\small
\renewcommand{\arraystretch}{1.3}
\begin{tabular}{@{}>{\raggedright\arraybackslash}p{3.8cm}
                   >{\raggedright\arraybackslash}p{4.1cm}
                   >{\raggedright\arraybackslash}p{4.6cm}@{}}
\toprule
Translation & Effect on the value & Price\\
\midrule
verifier $\to$ CSP & $\val(\Phi_V)=\max_\pi\Pr[V^\pi\text{ acc.}]$ exactly & none; $2^r$ constraints\\
CSP $\to$ verifier & the same, exactly & none (pick a random constraint)\\
constraint graph $\to$ game (double cover, Prop.~\ref{prop:cover}) & $\val(G)\le\val(G^{\times2})\le2\val(G)$ & soundness at most doubles\\
game $\to$ verifier (answer tables, Lemma~\ref{lem:tables}) & exactly the value of the game & $2\lceil\log\sigma\rceil$ binary queries\\
$q$-ary CSP $\to$ projection game (Thm.~\ref{thm:cv}) & $1-\frac{1-\val(\Phi)}q\ge\val(G_\Phi)\ge\val(\Phi)$ & gap shrinks by a factor $q$; alphabet $|\Sigma|^q$\\
verifier $\to$ clique (Thm.~\ref{thm:fglss}) & $\omega(H_x)=2^r\max_\pi\Pr[V^\pi\text{ acc.}]$ & $2^{r+q}$ vertices\\
\bottomrule
\end{tabular}
\caption{The price of every translation in the dictionary of Table~\ref{tab:dict}.}
\label{tab:price}
\end{table}


%% ===========================================================================
\section{The ladder of constraint types}
\label{sec:ladder}
%% ===========================================================================

Why are we so careful about the type of a constraint?  Because it determines both what an algorithm can do and what a hardness reduction needs.
We saw in Proposition~\ref{prop:ug-easy} that the \emph{unique} type makes perfect completeness easy: a label on one vertex forces the label on its neighbours in both
directions, so a satisfying labeling is found by propagation along a spanning tree.  For a projection the forcing goes in one direction only, and for a general relation in neither.
A $2$-to-$2$ constraint sits exactly between: after one step of propagation \emph{two} candidates remain, after $t$ steps up to $2^t$.

This is the right way to read the hierarchy of conjectures and theorems that organizes the rest of the course.  Every row of Table~\ref{tab:ladder} is a
statement about the complexity of a gap problem for constraint graphs of this type.

\begin{table}[t]
\centering\small
\renewcommand{\arraystretch}{1.3}
\begin{tabular}{@{}>{\raggedright\arraybackslash}p{2.6cm}
                   >{\raggedright\arraybackslash}p{3.6cm}
                   >{\raggedright\arraybackslash}p{3.3cm}
                   >{\raggedright\arraybackslash}p{2.9cm}@{}}
\toprule
Constraints & Gap problem & Status & Main use\\
\midrule
projection (large alphabet) & $\Gap{1}{\delta}{LabelCover}$, every $\delta>0$ & theorem (Thm.~\ref{thm:lc}) & H\aa stad's tight $3$-bit results (Lecture~\ref{lec:hastad})\\
$2$-to-$2$ & $\Gap{1-\eps}{\eps}{2-to-2}_k$ for some $k=k(\eps)$ & theorem \citep{KMS17,DKKMS18,KMS18} & \sat{Vertex Cover} within $\sqrt2-\eps$ (Lectures~13--15)\\
unique & $\Gap{1-\eps}{\delta}{UG}_k$ for some $k=k(\eps,\delta)$ & conjecture \citep{Khot02}; claimed in \citep{OpenAI26}, not verified & tight bounds for \sat{Max-Cut}, \sat{Vertex Cover}, every CSP (Lectures~8--11)\\
unique, completeness $1$ & $\Gap{1}{\delta}{UG}$ & in $\Ptime$ (Prop.~\ref{prop:ug-easy}) & ---\\
\bottomrule
\end{tabular}
\caption{The ladder of constraint types.  Moving down the table the constraints become more rigid and the hardness
results become stronger \emph{per label}, but harder to prove.}
\label{tab:ladder}
\end{table}

The $2$-to-$2$ theorem is called ``half-way'' to the UGC for a reason that can be proved in a few lines, and the proof is a good test of the language we have built.

\begin{lemma}[A $2$-to-$2$ relation is two bijections]\label{lem:split}
Let $C\subseteq\Sigma\times\Sigma'$ be a relation in which every label of $\Sigma$ and every label of $\Sigma'$ has exactly two partners; in particular
$C$ may be a $2$-to-$2$ constraint in the sense of Definition~\ref{def:types}.  Then $C$ is the disjoint union $M\sqcup M'$ of two perfect matchings, \ie of the graphs of two bijections.
\end{lemma}
\begin{proof}
The bipartite graph $(\Sigma\sqcup\Sigma',C)$ is simple and $2$-regular, hence a disjoint union of cycles, each of even length at least~$4$.  Colour the edges of each cycle alternately red and blue.
Every vertex lies on exactly one red and one blue edge, so the red edges form a perfect matching $M$ and the blue edges another one, $M'$.
\end{proof}

\begin{corollary}[Splitting a $2$-to-$2$ game]\label{cor:half}
Let $G$ be a constraint graph whose relations are all as in Lemma~\ref{lem:split}.  Let $\mathcal U_G$ be the Unique Games instance on the same vertices
in which every edge $e$ of $G$ is replaced by two parallel edges, one carrying each of the two bijections of $C_e$.  Then for every labeling $\ell$, the fraction of edges of $\mathcal U_G$
satisfied by $\ell$ is exactly \emph{half} the fraction of edges of $G$ satisfied by $\ell$.  Consequently $\val(\mathcal U_G)=\tfrac12\val(G)$.
\end{corollary}
\begin{proof}
A labeling $\ell$ satisfies $e=(u,v)$ in $G$ iff $(\ell(u),\ell(v))\in C_e=M_e\sqcup M'_e$, \ie iff it satisfies exactly one of the two parallel edges of $\mathcal U_G$ (the matchings
are disjoint, so never both); otherwise it satisfies neither.  So $\ell$ satisfies as many edges of $\mathcal U_G$ as of $G$, while $\mathcal U_G$ has twice as many edges.
\end{proof}

\noindent
Parallel edges are harmless: they only change the weights.  Hence the $2$-to-$2$ theorem gives the following unconditional
statement, which differs from the Unique Games Conjecture only in the \emph{completeness}.

\begin{corollary}[A half-way Unique Games Conjecture, from Table~\ref{tab:ladder}]\label{cor:halfway}
For every $\eps>0$ there is a $k=k(\eps)$ such that it is $\NP$-hard to distinguish Unique Games instances with $k$ labels and $\val\ge\tfrac{1-\eps}2$ from those with $\val\le\tfrac\eps2$.
\end{corollary}

\noindent
This is the sense in which the $2$-to-$2$ theorem proves the UGC ``half-way'' \citep{Khot19}: completeness $\approx\tfrac12$ instead of $1-\eps$.  Khot remarks that the arguments usually cited as evidence against the UGC apply to this
version too.  The claim of \citet{OpenAI26} is a claim to supply the missing half; we treat it as unverified (Lecture~16).

\begin{remark}[Max-3-Cut is $2$-regular]
The relation $x\ne y$ on three labels has exactly two partners for every label, so Lemma~\ref{lem:split} applies: it is the union of the two translations $y=x+1$ and $y=x+2$ modulo~$3$.
(It is not a union of $K_{2,2}$'s, since the bipartite graph is a $6$-cycle, so it is not $2$-to-$2$ in the sense of Definition~\ref{def:types}; see Exercise~\ref{ex:3cut}.)
\end{remark}


%% ===========================================================================
\section{From PCPs to cliques: the FGLSS graph}
\label{sec:fglss}
%% ===========================================================================

All of the above treated the proof as an oracle to be queried.  The reduction of \citet{FGLSS96} does something different: it turns the \emph{set of everything the verifier might see} into a
graph, and lets consistency do the work.  We fix a verifier $V$ as in Definition~\ref{def:pcp}, and an input $x$.  For $\rho\in\{0,1\}^r$ write $I_\rho=(i_1(\rho),\dots,i_q(\rho))$ for the positions it queries (possibly with repetitions).

\begin{definition}[FGLSS graph]\label{def:fglss}
A \defn{local view} is a pair $(\rho,a)$ with $\rho\in\{0,1\}^r$ and $a\in\{0,1\}^q$ such that $P_{x,\rho}(a)=1$ and $a_j=a_{j'}$ whenever $i_j(\rho)=i_{j'}(\rho)$.
Two local views $(\rho,a)$ and $(\rho',a')$ are \defn{compatible} if $\rho\neq\rho'$ and $a_j=a'_{j'}$ whenever $i_j(\rho)=i_{j'}(\rho')$.  The \defn{FGLSS graph} $H_x$ has the local views as vertices and an edge between
any two compatible views.  It has at most $2^{r+q}$ vertices.
\end{definition}

In words: a vertex says ``on random string $\rho$, the proof contained the accepting answers $a$''; two vertices are adjacent when they could both occur in one and the same proof.

\begin{theorem}[FGLSS]\label{thm:fglss}
$\displaystyle\omega(H_x)\;=\;2^r\cdot\max_{\pi\in\{0,1\}^N}\Pr_\rho\bigl[V^\pi(x)\text{ accepts}\bigr].$
\end{theorem}
\begin{proof}
($\ge$) Fix a proof $\pi$ and let $R=\{\rho:V^\pi(x)\text{ accepts on }\rho\}$.  For $\rho\in R$ the pair $(\rho,\pi|_{I_\rho})$ is a local view (accepting, and
consistent since it is read from the single string $\pi$).  Any two of these views agree with $\pi$ and have different $\rho$, so they are compatible: a clique of size $|R|=2^r\Pr_\rho[V^\pi\text{ accepts}]$.

($\le$) Let $K$ be a clique.  Adjacent vertices have different $\rho$, so $K$ contains at most one view per~$\rho$.  Define a proof $\pi$ by setting $\pi_i=a_j$ whenever $(\rho,a)\in K$ and $i_j(\rho)=i$, and $\pi_i=0$ for unconstrained positions.
This is well defined: two occurrences of position $i$ within one view agree by the consistency condition of a local view, and two occurrences in different views agree by compatibility.
Now for every $(\rho,a)\in K$ we have $\pi|_{I_\rho}=a$, so $P_{x,\rho}(\pi|_{I_\rho})=P_{x,\rho}(a)=1$: $V^\pi$ accepts on~$\rho$.  Hence $V^\pi$ accepts on at least $|K|$ random strings.
\end{proof}

\begin{figure}[t]
\centering
\begin{tikzpicture}[x=1cm,y=1cm,every node/.style={font=\small},
   vtx/.style={circle,draw=black!70,minimum size=0.78cm,inner sep=0pt,font=\footnotesize\ttfamily}]
  \node[vtx,fill=accent!22] (n001) at (111.00:3.00) {01};
  \node[vtx,fill=accent!22] (n010) at (135.00:3.00) {10};
  \node[vtx,fill=accent!22] (n011) at (159.00:3.00) {11};
  \node[vtx,fill=warn!22] (n100) at (21.00:3.00) {00};
  \node[vtx,fill=warn!22] (n101) at (45.00:3.00) {01};
  \node[vtx,fill=warn!22] (n111) at (69.00:3.00) {11};
  \node[vtx,fill=claim!18] (n200) at (-69.00:3.00) {00};
  \node[vtx,fill=claim!18] (n201) at (-45.00:3.00) {01};
  \node[vtx,fill=claim!18] (n211) at (-21.00:3.00) {11};
  \node[vtx,fill=black!14] (n300) at (-159.00:3.00) {00};
  \node[vtx,fill=black!14] (n301) at (-135.00:3.00) {01};
  \node[vtx,fill=black!14] (n310) at (-111.00:3.00) {10};
  \begin{scope}[on background layer]
    \draw[black!35] (n001)--(n101);
    \draw[black!35] (n001)--(n211);
    \draw[black!35] (n001)--(n310);
    \draw[black!35] (n010)--(n200);
    \draw[black!35] (n010)--(n201);
    \draw[black!35] (n010)--(n300);
    \draw[black!35] (n010)--(n301);
    \draw[black!35] (n011)--(n211);
    \draw[black!35] (n100)--(n200);
    \draw[black!35] (n100)--(n201);
    \draw[black!35] (n100)--(n300);
    \draw[black!35] (n100)--(n301);
    \draw[black!35] (n101)--(n211);
    \draw[black!35] (n101)--(n310);
    \draw[black!35] (n111)--(n211);
    \draw[black!35] (n200)--(n300);
    \draw[black!35] (n201)--(n301);
  \end{scope}
  \draw[accent!85!black,line width=1.8pt] (n011)--(n111);
  \draw[accent!85!black,line width=1.8pt] (n011)--(n310);
  \draw[accent!85!black,line width=1.8pt] (n111)--(n310);
  \node[vtx,fill=accent!22,draw=accent!85!black,line width=1.4pt] at (n011) {11};
  \node[vtx,fill=warn!22,draw=accent!85!black,line width=1.4pt] at (n111) {11};
  \node[vtx,fill=black!14,draw=accent!85!black,line width=1.4pt] at (n310) {10};
  \node[align=center,font=\footnotesize] at (135:4.55) {$\rho=1$: $x_1\vee x_2$\\ answers $x_1x_2$};
  \node[align=center,font=\footnotesize] at (45:4.55)  {$\rho=2$: $\neg x_1\vee x_2$\\ answers $x_1x_2$};
  \node[align=center,font=\footnotesize] at (-45:4.55) {$\rho=3$: $\neg x_2\vee x_3$\\ answers $x_2x_3$};
  \node[align=center,font=\footnotesize] at (-135:4.55){$\rho=4$: $\neg x_2\vee\neg x_3$\\ answers $x_2x_3$};
\end{tikzpicture}
\caption{The FGLSS graph of the verifier for $\varphi_0$ (Example~\ref{ex:phi0}) that picks one of the four clauses and reads its two variables.
Each group of three vertices contains the three satisfying answers of its clause; edges join compatible views.  The heavy triangle is a maximum clique: it consists of the views
of $\rho=1,2,4$ that agree with the assignment $x=(1,1,0)$, which satisfies $C_1,C_2,C_4$ and violates $C_3$.
No view of $\rho=3$ is compatible with all three.}
\label{fig:fglss}
\end{figure}

\begin{example}[The FGLSS graph of $\varphi_0$]\label{ex:fglss}
Take $V$ to be the verifier that expects an assignment $\pi\in\{0,1\}^3$ of $\varphi_0$ (Example~\ref{ex:phi0}), picks one of the four clauses ($r=2$) and reads its two variables ($q=2$).
Each clause has three satisfying answers, so $H_x$ has $4\cdot3=12$ vertices and (by direct count) $20$ edges; Figure~\ref{fig:fglss}.  Every assignment satisfies exactly three of
the four clauses, so $\max_\pi\Pr[\text{accept}]=\tfrac34$ and Theorem~\ref{thm:fglss} predicts $\omega(H_x)=4\cdot\tfrac34=3$, as one finds by inspection (there are $8$ maximum cliques, one for each
assignment).
\end{example}

\begin{pitfall}[unbreakable]{a clique is a consistent set, not just an accepting one}
Two vertices of $H_x$ may both be accepting views and still be non-adjacent: they disagree about one position of the proof.  A clique does not say that many random strings accept
\emph{somewhere}; it says that they accept \emph{the same proof}.  This is the same idea as the difference between the two provers of a game, who need not be consistent, and a single proof, which is:
Theorem~\ref{thm:fglss} is the one-oracle statement, and it is exact; no factor $2$ appears.
\end{pitfall}

\begin{corollary}[Constant-factor hardness for Clique]\label{cor:clique-const}
Unless $\Ptime=\NP$, no polynomial-time algorithm outputs, for every graph, a clique of more than half the maximum size.
\end{corollary}
\begin{proof}
By Theorem~\ref{thm:pcp1} there is a verifier for $\sat{3SAT}$ with $r=O(\log n)$, $q=O(1)$, perfect completeness and soundness $\tfrac12$.  Build $H_x$ in time polynomial in $n$ ($2^{r+q}=\poly(n)$ vertices).
If $x\in\sat{3SAT}$ then $\omega(H_x)=2^r$; otherwise $\omega(H_x)\le2^{r-1}$.  By Lemma~\ref{lem:gap-approx} no polynomial-time algorithm achieves ratio $\alpha>\tfrac12$ unless $\Ptime=\NP$.
\end{proof}

To pass from the constant $2$ to larger factors we amplify the gap.  Here graph products do exactly what independent repetition does for proofs.  The \defn{strong product} $G\boxtimes H$ has vertex set $V(G)\times V(H)$, and
two distinct vertices $(g,h)$, $(g',h')$ are adjacent iff ($g=g'$ or $g\sim g'$) and ($h=h'$ or $h\sim h'$).

\begin{proposition}[Clique number is multiplicative]\label{prop:product}
$\omega(G\boxtimes H)=\omega(G)\,\omega(H)$.
\end{proposition}
\begin{proof}
If $A\subseteq V(G)$ and $B\subseteq V(H)$ are cliques then $A\times B$ is a clique, giving $\ge$.  Conversely let $K$ be a clique of $G\boxtimes H$ and let $A$ and $B$ be its
projections to $V(G)$ and $V(H)$.  Two distinct elements of $A$ come from two vertices of $K$ that are adjacent, and since their first coordinates differ they must be adjacent in $G$; so $A$ is a clique, and likewise $B$.
Since $K\subseteq A\times B$ we get $|K|\le|A||B|\le\omega(G)\omega(H)$.
\end{proof}

\begin{corollary}[Hardness of Clique at larger factors]\label{cor:clique-quasi}
(a) For every constant $c$, it is $\NP$-hard to find a clique of size at least $1/c$ times the maximum.
(b) For every $\eps>0$, if on $N$-vertex graphs one can find in polynomial time a clique of size at least $2^{-(\log N)^{1-\eps}}$ times the maximum, then $\NP\subseteq\mathsf{DTIME}(2^{\polylog n})$.
\end{corollary}
\begin{proof}
Let $H=H_x$ be the graph of Corollary~\ref{cor:clique-const}, with $M=\poly(n)$ vertices, say $M\le n^{a}$.  The $k$-fold power $H^{\boxtimes k}$ has $M^k$ vertices and clique number $2^{rk}$ (YES) or at most $2^{(r-1)k}$ (NO), by Proposition~\ref{prop:product}: a gap of $2^k$.
For (a) take $k=\lceil\log_2c\rceil+1$, a constant.  For (b) take $k=\lceil(\log n)^{1/\eps}\rceil$: the graph has $N=n^{O(k)}$ vertices, is built in time $\poly(N)=2^{O((\log n)^{1+1/\eps})}$, and since
$\log N=\Theta(k\log n)$ the gap is $2^k=2^{\Theta((\log N)^{1/(1+\eps)})}$, which exceeds $2^{(\log N)^{1-\eps}}$ for large $n$ because $\tfrac1{1+\eps}>1-\eps$.  An algorithm as in (b) would therefore decide $\sat{3SAT}$ in quasi-polynomial time.
\end{proof}

\noindent
This is the shape of the argument of \citet{FGLSS96}, who had a PCP for $\NP$ with randomness and queries $O(\log n\log\log n)$ rather than $O(\log n)$ and $O(1)$; Theorem~\ref{thm:fglss}
applied to their proof system gives graphs with $n^{O(\log\log n)}$ vertices, and hence hardness under the assumption $\NP\not\subseteq\mathsf{DTIME}(n^{O(\log\log n)})$.  The PCP theorem removes the loss.

\begin{bigpicture}{Where the polynomial factors come from}
Corollary~\ref{cor:clique-quasi} reaches $2^{(\log N)^{1-\eps}}$ but not $N^{\Omega(1)}$, because $k$ independent repetitions cost $rk$ random bits, and $2^{rk}$ vertices is quasi-polynomial once $k$ grows with $n$.  The remedy is
\emph{randomness-efficient repetition}: replace the $k$ independent random strings by a $k$-step random walk on a constant-degree expander, which needs only $r+O(k)$ random bits and still drives the soundness down
exponentially in $k$ (Exercise~\ref{ex:expander}).  With $k=\Theta(\log n)$ one gets graphs of polynomial size and gap $N^{\Omega(1)}$, which is the conclusion of Theorem~6 of \citet{ALMSS98} (we do not claim that their proof proceeds this way).
The exponent was improved by the \emph{free bit} analysis of \citet{BGS98} and finally to $N^{1-\eps}$ by \citet{Hastad99}; the latter proof uses Label Cover, parallel repetition and the long code, the toolkit of Lectures~5 and~7.
\end{bigpicture}


%% ===========================================================================
\section{What Label Cover is for}
\label{sec:for}
%% ===========================================================================

The FGLSS reduction reads a PCP and builds a graph.  For problems with a rigid local structure---\sat{Max-E3-SAT}, \sat{Max-E3-Lin}, and later \sat{Max-Cut}---a different recipe is needed, and Label Cover is its first ingredient.

\begin{bigpicture}{The recipe of the next lectures}
\textbf{Step 1 (Lectures 1--2, 5).}  Prove that a Label Cover instance (projection constraints, soundness $\delta$, alphabet $\poly(1/\delta)$) is $\NP$-hard to approximate.
This is the \emph{outer verifier}.
\textbf{Step 2 (Lecture~\ref{lec:hastad}).}  Replace every vertex $v$ of the instance by a table $f_v$ that is supposed to be the \emph{long-code} encoding of the label of~$v$, and every edge $e=(u,v)$ by a
\emph{local test} between $f_u$ and $f_v$ that reads very few entries; this is the \emph{inner verifier}.  The tests have the shape of the target problem: three bits for \sat{Max-E3-Lin}, two bits for \sat{Max-Cut}.
\textbf{Step 3 (analysis).}  If the Label Cover instance has a labeling of value~$1$, encoding it passes all tests (completeness).  If a family of tables passes the tests with
probability above the claimed soundness, \emph{decode} each table into a short list of candidate labels; a random choice from the lists is a labeling that satisfies many edges, contradicting the soundness of
Label Cover.  The decoding is where the \emph{type} of the constraints matters: a projection lets the table at $v$ be read from the table at $u$, and the \emph{unique} type of the UGC lets decoding lose nothing
when the two tables are asked only two bits.
\end{bigpicture}

Label Cover is also the source of hardness for problems that look nothing like constraint satisfaction.  The same two-step recipe gives the logarithmic hardness of \sat{Set Cover}
\citep{LundYannakakis94,Feige98,DinurSteurer14}; and Label Cover under its original name was introduced in the context of lattice and code problems \citep{ABSS97}.


%% ===========================================================================
\section*{Exercises}
\addcontentsline{toc}{section}{Exercises}
%% ===========================================================================

\begin{exercise}[\stars{$\star$}]\label{ex:holonomy}
Let $G$ be a connected graph with linear constraints $x_v-x_u=c_e\pmod k$ on the edges $e=(u,v)$.  Prove that $\val(G)=1$ iff around every cycle the signed sum of the $c_e$ is $0$ mod $k$ (an edge traversed against
its orientation contributes $-c_e$).  \emph{Hint: spanning tree.}  Use this to give another proof of Proposition~\ref{prop:ug-easy} for linear constraints.
\end{exercise}

\begin{exercise}[\stars{$\star\star$}]\label{ex:cover-tight}
Compute $\val(G)$ and $\val(G^{\times2})$ for \sat{Max-Cut} on $K_n$ and show $\val(G^{\times2})/\val(G)=2-\tfrac2n$ for even $n$.  Show that for \sat{Max-Cut} on an odd cycle $C_\ell$ the ratio is $\tfrac\ell{\ell-1}$.  Conclude that the factor $2$ in Proposition~\ref{prop:cover} cannot be replaced by any constant smaller than~$2$.
\end{exercise}

\begin{exercise}[\stars{$\star\star$}]\label{ex:cv-tight}
(a) Let $\varphi$ be the $3$-CNF consisting of all eight clauses on three variables (every sign pattern).  Show that $\val(\varphi)=\tfrac78$ and $\val(G_\varphi)=\tfrac{23}{24}=1-\tfrac{1/8}3$, so the bound of Theorem~\ref{thm:cv} is tight.
(b) Show that the bound is also tight for \sat{Max-Cut} on a triangle, viewed as a CSP of arity $2$: $\val(\Phi)=\tfrac23$ and $\val(G_\Phi)=\tfrac56$.
\end{exercise}

\begin{exercise}[\stars{$\star\star$}]\label{ex:dreg}
Show that a relation $C\subseteq\Sigma\times\Sigma'$ in which every label of both sides has exactly $d$ partners is a disjoint union of $d$ perfect matchings.  \emph{Hint: Hall's theorem; remove one perfect matching and induct.}
Deduce the analogue of Corollary~\ref{cor:half}: from a $d$-to-$d$ game with $\val=1-\eps$ one gets a Unique Game with $\val=(1-\eps)/d$ on the same graph, with every edge repeated $d$ times.
\end{exercise}

\begin{exercise}[\stars{$\star$}]\label{ex:3cut}
(a) Verify that $x\ne y$ on $\Z_3$ is $2$-regular but not a disjoint union of copies of $K_{2,2}$.  (b) Write it as the union of the Unique constraints $y=x+1$ and $y=x+2$.  (c) For \sat{Max-3-Cut} on $K_4$ compute $\val(G)$ and the value of
the split Unique Games instance of Corollary~\ref{cor:half}.
\end{exercise}

\begin{exercise}[\stars{$\star\star\star$}]\label{ex:expander}
\emph{Randomness-efficient repetition.}  Let $A$ be the normalized adjacency matrix of a $d$-regular graph on the vertex set $\{0,1\}^r$ with second largest eigenvalue in absolute value at most $\lambda$, and let $B\subseteq\{0,1\}^r$ have density $\mu$.
(a) Prove that $\|P_BAP_B\|\le\mu+\lambda$, where $P_B$ is the coordinate projection onto $B$.  \emph{Hint: write a vector supported on $B$ as $\alpha\mathbf1/\sqrt{2^r}+y$ with $y\perp\mathbf1$ and bound $\alpha^2\le\mu\|u\|^2$ by Cauchy--Schwarz.}
(b) Deduce that a uniformly started random walk $\rho_1,\dots,\rho_k$ stays inside $B$ throughout with probability at most $\mu(\mu+\lambda)^{k-1}$.
(c) Let $V$ be a verifier with perfect completeness and soundness $\tfrac12$ and suppose an explicit family of such expanders with $\lambda\le\tfrac14$ and constant degree~$d$ exists.  Define $V_k$ that reads the walk's $k$ random strings, makes all $kq$ queries and
accepts iff all accept.  Show that $V_k$ uses $r+(k-1)\log d$ random bits and has soundness $(3/4)^{k-1}$, and conclude that with $k=\Theta(\log n)$ the FGLSS graph has $N=\poly(n)$ vertices and a gap $N^{\Omega(1)}$.
\end{exercise}

\begin{exercise}[\stars{$\star\star$}]\label{ex:reading2}
\emph{Reading.}  Section~2.3 of \citet{DIMACS10} reduces the PCP theorem to Label Cover by putting the random strings of the verifier on the left and the proof positions on the right.  Compare this reduction with Corollary~\ref{cor:verifier-game}.
Apply both to the verifier that picks a random clause of a $3$-CNF and queries its three variables, and identify the games.
\end{exercise}


%% ===========================================================================
\section*{Notes and further reading}
\addcontentsline{toc}{section}{Notes and further reading}
%% ===========================================================================

\paragraph{Where this lecture comes from.}
The dictionary between verifiers, CSPs and games is standard; we followed the treatment of the PCP theorem as a statement about Label Cover in Section~2.3 of the DIMACS notes \citep{DIMACS10}.  The name \emph{Label Cover}
comes from \citet{ABSS97}; two-prover one-round proof systems are due to \citet{FeigeLovasz92}; the clause--variable game is the canonical example.  The $d$-to-$1$ constraints are from \citet{Khot02}, and the $2$-to-$2$ constraints and the
``half-way'' remark are from \citet{Khot19}; the passage from a Unique Game to the bipartite two-prover game, with a loss of a factor of at most~$4$, is noted in \citet{Khot10}; we sharpened the factor to~$2$ for
the symmetric double cover (Proposition~\ref{prop:cover}).

\paragraph{On attributions.}
The Clique reduction and the exact formula $\omega(H_x)=2^r\cdot\max_\pi\Pr[\text{accept}]$ are due to \citet{FGLSS96}.  The hardness of approximating Clique within $N^\eps$ for some $\eps>0$, assuming only $\Ptime\ne\NP$, is Theorem~6 of
\citet{ALMSS98}, improving the $2^{\sqrt{\log N}}$ of \citet{AroraSafra98}; the factor $N^{1-\eps}$ is \citet{Hastad99}, as mentioned in Lecture~\ref{lec:why}.  The multiplicativity of the clique number under the strong
product is folklore.

\paragraph{What is not covered.}
We have not proved any hardness result for Label Cover itself beyond what Lecture~\ref{lec:why} did, and we have not discussed the \emph{size} of Label Cover instances (the Projection Games Conjecture asks for a nearly linear size).
Raz's theorem, which supplies the soundness, is quoted here; Lecture~\ref{lec:raz} proves the special case that Label Cover needs (expanding projection games) and states the general theorem.
